\documentclass[10pt,a4paper]{amsart}
\usepackage[margin=19mm]{geometry}
\usepackage{amsmath,amssymb,mathtools}
\usepackage[scr=rsfs,cal=euler]{mathalfa}
\usepackage{xfrac}
\usepackage[unicode,hidelinks,bookmarks=false]{hyperref}
\newcommand{\ie}{i.e.\ }
\newcommand{\cf}{cf.\ }
\newcommand{\resp}{respectively\ }
\newcommand{\Subsection}{\subsection}
\providecommand{\nobreakdash}{\nobreak\hbox{-}}
\setlength{\emergencystretch}{2em}
\multlinegap0pt
\usepackage[shortlabels]{enumitem}
\usepackage{amssymb}
\usepackage{mathtools}
\usepackage{dsfont}
\usepackage{bm}
\usepackage{algorithm}\usepackage{algpseudocode}
\usepackage{tikz}
\usepackage{xcolor}
\newtheorem{theorem}{Theorem}
\newcommand{\Let}{\triangleq}
\DeclareMathOperator*{\argmax}{arg\,max}
\DeclareMathOperator*{\argmin}{arg\,min}
\newcommand{\opt}{^\star}
\title{Selective Ambulance Dispatch Under Contextual Travel-Time Uncertainty}
\author{Zikun Lin, Daniel Zhuoyu Long, Viet Anh Nguyen}
\date{Source: arXiv:2605.23378v1 (CC BY 4.0)}
\begin{document}
\maketitle

\noindent\emph{Source and licence.} Selective Ambulance Dispatch Under Contextual Travel-Time Uncertainty. Version arXiv:2605.23378v1 (CC BY 4.0), distributed by arXiv under the Creative Commons Attribution 4.0 International licence, http://creativecommons.org/licenses/by/4.0/, as recorded in the arXiv licence metadata for this exact version and shown on the abstract page https://arxiv.org/abs/2605.23378v1. Full licence notice: \texttt{licenses/arxiv-2605.23378-CC-BY-4.0.txt}.\par\smallskip

\noindent\textit{Editorial scope.} Counted unit: the article's theorem thm:ideal\_optimality (source0.tex lines 358--371). The article's complete original proof of it is delivered with it (source0.tex lines 1219--1300, one \texttt{proof} environment of 82 lines, which the article places away from the statement). No mathematics is written, repaired or re-derived: the statement and the proof are the article's own lines, in the article's own notation, and no retained line is substituted or re-flowed.  Retained uncounted as the source-local context the proof needs: lines 300--305; lines 319--324.  Nothing was omitted from any retained span, so the package carries no omission record.  No retained line contains a citation, so the wrapper carries no bibliography and no bibliography entry.  No retained line contains a figure environment, an image-inclusion command or drawing code, so no image is delivered and the package declares no asset dependency.  Editorial formatting only, in addition: an AMS \texttt{amsart} preamble that restates the article's own declarations and macros used by the retained lines, namely the environments \texttt{theorem} on separate counters, together with the standard packages those lines need.  The retained environments print in this document as Theorem 1 in order of appearance, whereas the article prints the same items with the numbering of its own full document.  The complete native source of the article as distributed in the arXiv e-print for this exact version is bundled unchanged as \texttt{sources/201203/source0.tex}.
\begin{equation}
\label{eq:delta-rho}
\delta(\mathcal T; z_1) \Let
\max_{c\in\mathcal C(\mathcal T)}
\Bigl( c^\top z_1  -  \underbrace{\min_{z \in\mathcal Z} c^\top z}_{\text{fastest arrival under scenario } c} \Bigr).
\end{equation}

\begin{equation}
\label{eq:decision-rule}
\text{dispatch a second ambulance along } z_2
\quad\Longleftrightarrow\quad
\delta(\mathcal T; z_1) > \mathrm{thr}(\mathcal T).
\end{equation} 

\begin{theorem}[Optimality of the IDEAL policy]
\label{thm:ideal_optimality}
Fix a context $\mathcal T$ and a primary path $z_1 \in \mathcal Z$. Consider the optimization problem
\begin{equation}
\label{eq:pthr}
\tag{$\mathcal P_{\text{thr}}$}
\max_{\tau\in\{0,1\},\,z\in\mathcal Z} \Bigl\{ \tau\bigl( \Delta(z;\mathcal T,z_1) - \mathrm{thr}(\mathcal T) \bigr) \Bigr\}.
\end{equation}
When $\mathcal C(\mathcal T)$ is nonempty and compact, the IDEAL policy produces a pair $(\tau,z_2)$ that solves~\eqref{eq:pthr}, and the following two statements hold:
\begin{enumerate}[leftmargin=7mm]
\item The scenario $c\opt$ and secondary path $z_2$ constructed in the algorithm achieve the optimistic time gap $\Delta(z_2;\mathcal T,z_1) = \delta(\mathcal T; z_1)$.
\item The policy dispatches a second ambulance if and only if every optimal solution of~\eqref{eq:pthr} has $\tau=1$, which is equivalent to condition~\eqref{eq:decision-rule}.
\end{enumerate}
\end{theorem}

\begin{proof}[Proof of Theorem~\ref{thm:ideal_optimality}]
Fix a context $\mathcal T$ and a primary path $z_1 \in \mathcal Z$. Let $\delta(\mathcal T; z_1)$ be the optimistic time gap defined in~\eqref{eq:delta-rho}. Since $\mathcal Z$ is finite, the map $c\mapsto \min_{z\in\mathcal Z}c^\top z$ is the pointwise minimum of finitely many linear functions and is continuous. Hence $c\mapsto c^\top z_1-\min_{z\in\mathcal Z}c^\top z$ is continuous. Since $\mathcal C(\mathcal T)$ is nonempty and compact, the maximum in~\eqref{eq:delta-rho} is attained. Hence there exists
\[
c\opt \in \argmax_{c\in\mathcal C(\mathcal T)}
\bigl(c^\top z_1 - \min_{z\in\mathcal Z} c^\top z\bigr),
\]
and we can choose an associated secondary path $z_2 \in \argmin_{z\in\mathcal Z} (c\opt)^\top z$.

We first show that the pair $(c\opt,z_2)$ constructed by the IDEAL algorithm achieves the optimistic time gap in the sense that
\begin{equation}
\label{eq:delta-equals-Delta-max}
\Delta(z_2;\mathcal T,z_1) = \delta(\mathcal T; z_1)
= \max_{z\in\mathcal Z} \Delta(z;\mathcal T,z_1),
\end{equation}
and then show that the pair produced by IDEAL is an optimal solution of~\eqref{eq:pthr}, with the convention that zero surplus is resolved by choosing $\tau=0$.

For any secondary path $z\in\mathcal Z$ and any scenario $c\in\mathcal C(\mathcal T)$,
$c^\top z_1 - c^\top z \le c^\top z_1 - \min_{z\in\mathcal Z} c^\top z$. Taking the maximum over $c\in\mathcal C(\mathcal T)$ on both sides yields
\[
\Delta(z;\mathcal T,z_1)
= \max_{c\in\mathcal C(\mathcal T)} \bigl(c^\top z_1 - c^\top z\bigr)
\le
\max_{c\in\mathcal C(\mathcal T)}\bigl(c^\top z_1 - \min_{z\in\mathcal Z} c^\top z\bigr)
= \delta(\mathcal T; z_1),
\]
for every $z\in\mathcal Z$. Thus $\delta(\mathcal T; z_1)$ is an upper bound on $\Delta(z;\mathcal T,z_1)$ for all secondary paths $z$.

Now consider the specific path $z_2$ chosen by the algorithm. By construction, $(c\opt)^\top z_2 = \min_{z\in\mathcal Z} (c\opt)^\top z$, so
$(c\opt)^\top z_1 - (c\opt)^\top z_2
= (c\opt)^\top z_1 - \min_{z\in\mathcal Z} (c\opt)^\top z
= \delta(\mathcal T; z_1)$.
Therefore
\[
\Delta(z_2;\mathcal T,z_1)
= \max_{c\in\mathcal C(\mathcal T)} \bigl(c^\top z_1 - c^\top z_2\bigr)
\ge
(c\opt)^\top z_1 - (c\opt)^\top z_2
= \delta(\mathcal T; z_1).
\]
Combining this inequality with the upper bound $\Delta(z_2;\mathcal T,z_1)\le \delta(\mathcal T; z_1)$ gives
\[
\Delta(z_2;\mathcal T,z_1) = \delta(\mathcal T; z_1)
\quad\text{and}\quad
\Delta(z_2;\mathcal T,z_1) = \max_{z\in\mathcal Z} \Delta(z;\mathcal T,z_1),
\]
which proves~\eqref{eq:delta-equals-Delta-max} and establishes Item~1 in the theorem.

The optimization problem~\eqref{eq:pthr} can be analyzed by first fixing $z\in\mathcal Z$ and optimizing over $\tau\in\{0,1\}$. For a given $z$, the objective in~\eqref{eq:pthr} equals
\[
\max_{\tau\in\{0,1\}} \tau\bigl(\Delta(z;\mathcal T,z_1) - \mathrm{thr}(\mathcal T)\bigr)
=
\begin{cases}
0, & \text{if } \Delta(z;\mathcal T,z_1) \le \mathrm{thr}(\mathcal T),\\[0.5ex]
\Delta(z;\mathcal T,z_1) - \mathrm{thr}(\mathcal T) & \text{if } \Delta(z;\mathcal T,z_1) > \mathrm{thr}(\mathcal T).
\end{cases}
\]
Hence~\eqref{eq:pthr} is equivalent to
\[
\max_{z\in\mathcal Z}
\max\bigl\{0, \Delta(z;\mathcal T,z_1) - \mathrm{thr}(\mathcal T)\bigr\}.
\]
Let $\Delta_{\max} \Let \max_{z\in\mathcal Z} \Delta(z;\mathcal T,z_1)$. Then the optimal value of~\eqref{eq:pthr} equals $\max\{0,\Delta_{\max} - \mathrm{thr}(\mathcal T)\}$, and any maximizer $z\opt$ of $\Delta(z;\mathcal T,z_1)$ is optimal for~\eqref{eq:pthr} when paired with
\[
\tau\opt =
\begin{cases}
1, & \text{if } \Delta_{\max} > \mathrm{thr}(\mathcal T),\\
0, & \text{if } \Delta_{\max} \le \mathrm{thr}(\mathcal T).
\end{cases}
\]
Moreover, every optimal solution of~\eqref{eq:pthr} has $\tau=1$ exactly when $\Delta_{\max} > \mathrm{thr}(\mathcal T)$. If $\Delta_{\max}\le \mathrm{thr}(\mathcal T)$, the choice $\tau=0$ attains the optimal value $0$.

The identity~\eqref{eq:delta-equals-Delta-max} shows that $\Delta_{\max} = \delta(\mathcal T; z_1)$ and that the secondary path $z_2$ constructed by the IDEAL algorithm is a maximizer of $\Delta(z;\mathcal T,z_1)$. The algorithm uses the rule
\[
\tau =
\begin{cases}
1, & \text{if } \delta(\mathcal T; z_1) > \mathrm{thr}(\mathcal T),\\
0, & \text{if } \delta(\mathcal T; z_1) \le \mathrm{thr}(\mathcal T),
\end{cases}
\]
which matches the optimal choice of $\tau\opt$ above. Thus the pair $(\tau,z_2)$ produced by the IDEAL policy solves~\eqref{eq:pthr}, and the policy dispatches a second unit if and only if $\delta(\mathcal T; z_1) > \mathrm{thr}(\mathcal T)$, 
which is precisely the decision rule~\eqref{eq:decision-rule}. This establishes Item~2 and completes the proof.
\end{proof}

\end{document}
